\documentclass[10pt]{article}
\usepackage[margin=20mm]{geometry}
\usepackage{amsmath,amssymb,mathtools,fontspec,unicode-math,xurl,hyperref}
\IfFontExistsTF{DejaVu Serif}{\setmainfont{DejaVu Serif}}{\setmainfont{Latin Modern Roman}}
\setmathfont{latinmodern-math.otf}
\setlength{\parindent}{0pt}\setlength{\parskip}{6pt}\sloppy
\title{Scientific review of
Brownian motion with asymptotically normal reflection in unbounded domains: From transience to stability}
\author{Initial scientific review: M25-0485}
\date{29 October 2026}
\begin{document}
\maketitle
Abstract: The acquired accepted manuscript contains an over-uniform probability bound in Proposition 7.2 and invalid steps in the divergent-integral calculation (62) used to prove the stationary-tail lower bound. Constant test functions refute the former as quantified. The latter calculation needs a different integration interval and preservation of a positive scaling factor. The assessment is Major for the inspected core lower-tail derivation. Explicit local repairs are given, and no counterexample to the final recurrence, tail, or convergence exponents is established.\par
\section*{Source, scope, and principal claims}
The paper is by Miha Brešar, Aleksandar Mijatović, and Andrew Wade, associated with Annals of Probability 53 (2025), 175–222, DOI 10.1214/24-AOP1703 (https://doi.org/10.1214/24-AOP1703). The reviewed source is the Warwick accepted manuscript deposited in April 2024: a repository cover followed by 55 manuscript pages. References below use printed manuscript pages; physical PDF pages are one greater. The 48-page journal typesetting has not been compared with these findings. The public arXiv record instead identifies v2 dated 7 April 2023.\par
The domain is \(\mathcal D=\{(x,y):x\ge0,\ \|y\|\le b(x)\}\). Assumptions (D2), (C2), and (V2) include the regular origin, bounded uniformly elliptic covariance, uniformly inward bounded reflection, sublinear boundary growth, and uniform asymptotics. Set \[\beta_c=\frac{c_0\sigma_1^2}{s_0\sigma_2^2},\qquad
 m_c=\frac{1-\beta/\beta_c}{2},\qquad
 M_c=-\frac{1+\beta/\beta_c}{2}.\] Theorem 1.3 treats recurrence and transience, with stronger assumptions at \(\beta=\beta_c\). Theorem 1.4 gives return-time exponents up to arbitrary positive error. Theorem 1.5 assumes \(\beta<-\beta_c\) and asserts stationary-tail exponent \(-2M_c\) and total-variation exponent \(-M_c\), again with arbitrary error. The boundary case \(\beta=-\beta_c\) is expressly omitted, so its absence is not a defect.\par
The detailed review concerns the Lyapunov signs, the additive-functional estimate, and the lower-tail argument in Section 7. The localization proof of Feller continuity, irreducibility, and the external Harris-recurrence criteria have not been independently certified. This scope does not amount to a complete verification of the 55-page manuscript.\par
\section*{A probability bound with incorrect uniformity}
Equation reference: eq:original. Problem 1 (Proposition 7.2, pp. 47–48).  For fixed model parameters, \(p>1-\beta/\beta_c\), a bounded measurable set \(D\), and a starting point \(z\), the proposition first chooses positive constants \(\delta,C,r_0,c_1,c_2,\epsilon\). It then claims, for every continuous nondecreasing \(H\ge0\) that is not identically zero, \[
 \mathbb P_z\!\left(\int_0^{\tau_\delta(D)}
 H(c_1(X_s+c_2))\,ds\ge r\right)
 \ge \frac{C}{G_H(r/\epsilon)^p},\qquad r>r_0,\] where \(G_H\) inverts \(v\mapsto v^2H(v)\). This order of quantifiers is impossible.\par
Exact contradiction. Let \(H_K(v)=K>0\). It meets every condition imposed on \(H\), with \(r_H=0\), and \(G_{H_K}(u)=\sqrt{u/K}\). For any fixed \(r>r_0\), the right side of [eq:original] (\#eq:original) is \[C\left(\frac{\epsilon K}{r}\right)^{p/2}.\] Taking \(K=(r/\epsilon)(2/C)^{2/p}\) makes it \(2\), whereas the left side is a probability and cannot exceed \(1\).\par
The model assumptions are nonempty even in the stable regime. For example, take \(d=1\), identity covariance, normal reflection, and \[b(x)=\frac{\sqrt{x}}{(1+x)^{5/2}}\quad(x>0),\qquad b(0)=0.\] Then \(bb'\to1/2\), \(b''/(b')^3\to-2\) at zero, and \(xb'/b\to-2\), \(b',b''\to0\) at infinity. The boundary is smooth at the tip: the equation \(y^2=x/(1+x)^5\) solves smoothly for \(x\) near \(0\). Normal reflection is a bounded smooth unit field with inward component \(1\); its asymptotic constants are \(s_0=c_0=1\). Thus \(\beta=-2\), \(\beta_c=1\), and \(p=4\) is admissible. The contradiction does not depend on computing this diffusion’s return distribution. ◻\par
\subsection*{Location of the loss and the necessary correction}
Equation (58) starts with a valid type of bound truncated by a minimum with \(1\), under a lower bound \(v>v_*\). Substituting \(v=G_H(r/\epsilon)\) requires \[r>\epsilon v_*^2 H(v_*).\] This threshold depends on \(H\). In equation (59), deleting the minimum additionally requires the inverse to be sufficiently large. A sufficient corrected form allows \(r_0=r_0(H)\), with the other model constants fixed, and asserts [eq:original] (\#eq:original) only beyond that threshold. Alternatively the truncated bound and its original domain can be retained.\par
The uniform assertion is false, but the later divergent-integral criterion uses a fixed \(H\); changing a finite lower integration limit does not alter divergence. This quantifier defect therefore has a local repair and no demonstrated independent effect on the final exponents. It is assessed as Minor.\par
\section*{The core divergent-integral calculation}
Problem 2 (Lemma 7.5, equation (62), pp. 49–50).  The lower-tail proof constructs a nondecreasing function \(H\), integrable under a hypothetical stationary measure with too-small tails, and seeks a contradictory divergent integral. The equalities and inequalities printed in (62) do not follow from that construction.\par
Write \(\eta>0\) for the error denoted by \(\epsilon\) in the construction, and set \[\alpha=-(1+\beta/\beta_c)+\eta,\qquad s=\alpha-\eta/2>0.\] For \(r_{n+1}>2r_n>2\), the construction has \(H=1\) before \(r_1\), derivative \(r_n^s\) on \([r_n,r_n+1)\), and derivative \((r_{n+1}-r_n)^{-1}\) on the rest of \([r_n,r_{n+1})\). It naturally defines a continuous locally absolutely continuous function; differentiability at the corners is unnecessary for the integral argument.\par
First, \(H(r_n)\) need not equal \(r_n^s\); the increment of that size is only accumulated by \(r_n+1\). For example, with \(s=3,r_1=10,r_2=10000\), the construction gives \[H(r_2)=1001+\frac{9989}{9990}<1002,\qquad r_2^s=10^{12}.\] This is a diagnostic of the claimed algebra, not a proposed stationary law. Arbitrarily rapid subsequences are permitted by the contradiction construction, so the growth condition \(r_{n+1}>2r_n\) cannot supply the missing lower bound at the left endpoints.\par
Second, replacing the integral of the increasing transformed function by its value at the left endpoint gives an inequality, not the printed equality. Third, integrating over endpoints \(c_1(r_n+c_2)\), \(c_1(2r_n+c_2)\) produces plus signs in the shifted powers, whereas the following line has minus signs. Finally, the display drops \(c_1^{2-p}\) in a lower bound even though \(c_1>1,p>2\), so this factor is less than \(1\). For any positive finite partial-sum term, that deletion increases the expression. These are failures of the printed calculation; the fact that a correctly formed infinite series diverges does not justify the finite algebra.\par
\subsection*{A verified local replacement}
The required divergence can nevertheless be obtained. Put \(\widetilde H(u)=H(u/c_1-c_2)\) beyond its constant initial extension, and let \(G\) invert \(u^2\widetilde H(u)\). Choose \[1-\beta/\beta_c<p<1-\beta/\beta_c+\eta/2 .\] Then \(p>2\) and \(s+2-p>0\). On the disjoint intervals \[I_n=[c_1(r_n+1+c_2),\ c_1(2r_n+c_2)]\] we have \(\widetilde H(u)\ge H(r_n+1)\ge r_n^s\). For all sufficiently large \(n\), the length of \(I_n\) is at least \(c_1r_n/2\), and \(u\le3c_1r_n\). Since \(1-p<0\), \[\int_{I_n}u^{1-p}\widetilde H(u)\,du
 \ge \frac{3^{1-p}c_1^{2-p}}{2}\,r_n^{s+2-p}.\] The coefficient is positive and is retained. Substitution \(r=\epsilon u^2\widetilde H(u)\), using the almost-everywhere derivative, now yields \[\int^\infty G(r/\epsilon)^{-p}\,dr
 =\epsilon\int^\infty u^{-p}
       (2u\widetilde H(u)+u^2\widetilde H'(u))\,du
 \ge 2\epsilon\sum_{n\ge n_0}
       \int_{I_n}u^{1-p}\widetilde H(u)\,du
 =\infty.\] The finiteness estimate (61) remains valid: its contributions are bounded by \(\sum_n r_n^{-\eta/2}+\sum_n r_n^{-\alpha}<\infty\). Thus the intended contradiction survives these explicit changes, conditional on the preceding stochastic and invariant-measure ingredients.\par
\subsection*{Dependency and impact}
The failed calculation is the decisive divergent-integral step of Lemma 7.5. That lemma supplies Proposition 7.4’s stationary-tail lower bound, which in turn feeds the total-variation lower bound and Theorem 1.5. This is a Major claim-level defect in the acquired derivation, with the dependency explicit. The local replacement above does not establish that the final theorem is false; it instead supplies substantial counterevidence to that stronger conclusion. It also does not certify the unexamined Feller and external recurrence arguments.\par
\section*{Other checks and limits}
The inspected Lyapunov formula has leading generator coefficient \[\gamma f_{w,1}^{\gamma-2}
 \{\sigma_1^2(\gamma-1)+\sigma_2^2(1-w)\},\] and reflection sign determined by \(\gamma(s_0\beta/c_0-1+w)\) (Lemma 2.1, pp. 13–14). The interval of \(w\) making both terms negative is nonempty precisely when \(0<\gamma<1-\beta/\beta_c\), in agreement with Lemma 5.4. For \(\gamma>2\), the drift-rate inversion in Proposition 7.1 gives exponent \(\gamma/2-1\), approaching \(M_c\). In Appendix A, the total-variation comparison \(\|\pi-P_t(z,\cdot)\|_{\rm TV}\ge f(y)-g(z,t)/y\) gives \(f(y)/2\) when \(yf(y)=2g(z,t)\), as required for the power functions actually used.\par
The heuristic Bessel dimension on p. 10 should be \(1+\beta/\beta_c\) after unit-variance normalization, since a Bessel generator has drift \((d_{\rm Bes}-1)/(2x)\). The printed \(2\beta/\beta_c\) is not used in the formal Lyapunov calculation; this is a clarification, not an additional substantive finding.\par
No conclusion is drawn about the omitted critical stability case, a journal-final correction, or the complete localization proof. Public title/correction searches and the arXiv and author records did not identify a separate erratum or supplement in the material retrieved. A journal issue-level search hit was not treated as a comparison of the final article.\par
Assessment | Repair verification | Author adoption | Remaining scope\par
Major: (62); Minor: uniformity in 7.2 | Fixed-\(H\) cutoff and divergent-integral replacement checked | Unknown | Journal comparison; full stochastic prerequisites\par
\section*{Scope and references}
The retained independent calculation uses symbolic limits for the explicit domain and exact rational finite diagnostics for the series steps. It does not execute author code or simulate the diffusion. Scientific reasoning and page inspection belong to the assigned initial-review AI session; the native helper performs only clerical binding checks. No human review or independent acceptance is asserted.\par
9 M. Brešar, A. Mijatović, and A. Wade, Brownian motion with asymptotically normal reflection in unbounded domains: from transience to stability. Warwick WRAP accepted manuscript, deposit April 2024, https://wrap.warwick.ac.uk/id/eprint/184881/. Journal DOI: 10.1214/24-AOP1703. The same authors, arXiv:2303.06916, version record (v2, 7 April 2023), consulted for version history only.\par
\end{document}
